% $Id$

\chapter{Introduction}

This is where we introduce PRoNTo and explain what we wanted to do...\\
Cite a few important references like \cite{Bishop2006}, or \cite{Geurts2006}, or another \cite{Mourao-Miranda2007} and obvious \cite{Rasmussen2006}.
\\ \\
This comes from the Harvest application document.

Advances in neuroimaging techniques have radically changed the way neuroscientists address questions about functional anatomy, especially in relation to behavioural and clinical disorders. Many questions about brain function, previously investigated using electrophysiological recordings in animals can now be addressed non-invasively in humans. Such studies have yielded important results in cognitive neuroscience and neuropsychology. Amongst the various neuroimaging modalities available, Magnetic Resonance Imaging (MRI) has become widely used due to its relatively high spatial and temporal resolution, and because it is safe and non-invasive. By selecting specific MRI sequence parameters, different MR signals can be obtained from different tissue types, giving images with high contrast among organs, between normal and abnormal tissues and/or between activated and deactivated brain areas. MRI is often sub-categorized into structural MRI (sMRI) and functional MRI (fMRI). Structural and functional MRI data are inherently multivariate in nature, since each scan contains information about tissue structure, or brain activation, at thousands of measured locations (voxels). Considering that most brain functions are distributed processes involving a network of brain regions, it would seem desirable to use the spatially distributed information contained in the data to give a better understanding of brain functions in normal and abnormal conditions. 

The typical analysis pipeline in neuroimaging is strongly rooted in a mass-univariate statistical approach, which assumes that activity in one brain region occurs independently from activity in other regions. Although this has yielded great insights over the years, and continues to be the tool of choice for data analysis, there is a growing recognition that the spatial dependencies among signal from different brain regions should be properly modeled. The effect of interest can be subtle and spatially distributed over the brain - a case of high-dimensional, multivariate data modeling for which conventional tools may lack sensitivity. 

Therefore, there has been an increasing interest in investigating this spatially distributed information using multivariate pattern recognition methods. Where pattern recognition has been used in neuroimaging, it has led to fundamental advances in the understanding of how the brain represents information and has been applied to many diagnostic applications. For the latter, this approach can be used to predict the status of the patient scanned (healthy vs. diseased or disease A vs. B) and can provide the discriminating pattern leading to this classification. 
Several active areas of research in machine learning are crucially important for the difficult problem of neuroimaging data analysis: modelling of high-dimensional multivariate time series, sparsity, regularisation, dimensionality reduction, causal modeling, and ensembling to name a few. However, the application of pattern recognition approaches to the analysis of neuroimaging data is limited mainly by the lack of “user-friendly” and comprehensive tools available to the fundamental, cognitive, and clinical neuroscience communities. Furthermore, it is not uncommon for these methods to be used incorrectly, with the most typical case being improper separation of training and testing datasets.

The aim of the project is to develop a toolbox based on machine learning techniques for the analysis of neuroimaging data (PRoNTo: Pattern Recognition Neuroimaging Toolbox). The toolbox will be a \matlab  based software package specifically designed for the analysis of brain imaging data including steps for feature extraction according to the experimental design, feature selection, machine training and testing. Initially it will focus on supervised learning approaches such as Support Vector Machines, Gaussian Processes, etc. Other techniques will be added once a consistent framework and data flow has been implemented.

The toolbox will facilitate the interaction between machine learning and neuroimaging communities. On one hand the machine learning community will be able to contribute to the toolbox with novel published machine learning algorithms. On the other hand the toolbox will provide a variety of tools for the neuroscience and clinical neuroscience communities, enabling them to ask new questions that cannot be easily investigated using existing statistical analysis tools. The toolbox code will be distributed for free, but as copyright software under the terms of the GNU General Public License as published by the Free Software Foundation.
